% TODO make hw stylesheet
\documentclass[10pt]{scrartcl}
\usepackage[a4paper, total={6in, 8in}]{geometry}
\usepackage[usenames,dvipsnames,svgnames]{xcolor}
\usepackage[shortlabels]{enumitem}
\usepackage[framemethod=TikZ]{mdframed}
\usepackage{amsmath,amssymb,amsthm}
\usepackage[colorlinks]{hyperref}
\usepackage{multicol}
\usepackage{tabularx}

\hypersetup{
	colorlinks=true,
	linkcolor=blue,
	filecolor=magenta,      
	urlcolor=magenta,
	pdftitle={MAT 314 HW}
}

\usepackage{microtype}
\usepackage{mathtools}
\usepackage[headsepline]{scrlayer-scrpage}
\usepackage{thmtools}
\usepackage[textsize=footnotesize]{todonotes}

\mdfdefinestyle{mdbluebox}{roundcorner=10pt,innerbottommargin=9pt,
	linecolor=blue,backgroundcolor=TealBlue!5,}
\declaretheoremstyle[headfont=\sffamily\bfseries\color{MidnightBlue},
mdframed={style=mdbluebox},]{thmbluebox}

\mdfdefinestyle{mdbloobox}{frametitlefont=\bfseries,innerbottommargin=8pt,
	nobreak=true,backgroundcolor=TealBlue!5,linecolor=blue,}
\declaretheoremstyle[headfont=\bfseries\color{MidnightBlue},
mdframed={style=mdbloobox},headpunct={\\[3pt]},postheadspace=0pt,]{thmbloobox}

\mdfdefinestyle{mdredbox}{frametitlefont=\bfseries,innerbottommargin=8pt,
	nobreak=true,backgroundcolor=Salmon!5,linecolor=RawSienna,}
\declaretheoremstyle[headfont=\bfseries\color{RawSienna},
mdframed={style=mdredbox},headpunct={\\[3pt]},postheadspace=0pt,]{thmredbox}

\mdfdefinestyle{mdgreenbox}{linecolor=ForestGreen,backgroundcolor=ForestGreen!5,
	linewidth=2pt,rightline=false,leftline=true,topline=false,bottomline=false,}
\declaretheoremstyle[headfont=\bfseries\sffamily\color{ForestGreen!70!black},
mdframed={style=mdgreenbox},headpunct={ --- },]{thmgreenbox}

\mdfdefinestyle{mdorangebox}{linecolor=RedOrange,backgroundcolor=RedOrange!5,
	linewidth=2pt,rightline=false,leftline=true,topline=false,bottomline=false,}
\declaretheoremstyle[headfont=\bfseries\sffamily\color{RedOrange!70!black},
mdframed={style=mdorangebox},headpunct={ --- },]{thmorangebox}

\mdfdefinestyle{mdblackbox}{linecolor=black,backgroundcolor=RedViolet!5!gray!5,
	linewidth=3pt,rightline=false,leftline=true,topline=false,bottomline=false,}
\declaretheoremstyle[mdframed={style=mdblackbox}]{thmblackbox}

\declaretheoremstyle[spaceabove=6pt,spacebelow=6pt]{basehead}

\declaretheorem[style=basehead,name=Problem]{problem}
\declaretheorem[style=basehead,name=Problem,numbered=no]{problem*}
\declaretheorem[style=thmbloobox,name=Setup]{setup} % for config geo
\declaretheorem[style=thmredbox,name=Example,sibling=problem]{example}
\declaretheorem[style=thmredbox,name=$\bigstar$ Problem,sibling=problem]{niceprob}
\declaretheorem[style=thmbluebox,name=Theorem,sibling=problem]{theorem}
\declaretheorem[style=thmbluebox,name=Corollary,sibling=theorem]{corollary}
\declaretheorem[style=thmbluebox,name=Proposition,sibling=theorem]{prop}
\declaretheorem[style=thmbluebox,name=Lemma,numberwithin=problem]{lemma}
\declaretheorem[style=thmbluebox,name=Theorem,numbered=no]{theorem*}
\declaretheorem[style=thmbluebox,name=Lemma,numbered=no]{lemma*}
\declaretheorem[style=thmgreenbox,name=Claim,numberwithin=problem]{claim}
\declaretheorem[style=thmgreenbox,name=Claim,numbered=no]{claim*}
\declaretheorem[style=thmblackbox,name=Remark,numberwithin=problem]{remark}
\declaretheorem[style=thmblackbox,name=Remark,numbered=no]{remark*}
\declaretheorem[style=thmgreenbox,name=Definition,sibling=theorem]{definition}
\declaretheorem[style=thmgreenbox,name=Definition,numbered=no]{definition*}
\declaretheorem[style=thmorangebox,name=Warning,sibling=theorem]{warning}
\declaretheorem[style=thmorangebox,name=Warning,numbered=no]{warning*}
\declaretheorem[style=thmorangebox,name=Note,numbered=no]{note}
\declaretheorem[style=thmblackbox,name=Exercise,sibling=problem]{exercise}
\declaretheorem[style=thmblackbox,name=Exercise,numbered=no]{exercise*}
\declaretheorem[style=thmblackbox,name=Question,sibling=problem]{question}
\declaretheorem[style=thmblackbox,name=Question,numbered=no]{question*}

\addtolength{\textheight}{3.14cm}
\providecommand{\ol}{\overline}
\providecommand{\ul}{\underline}
\providecommand{\eps}{\varepsilon}
\providecommand{\half}{\frac{1}{2}}
\providecommand{\dang}{\measuredangle} %% Directed angle
\providecommand{\CC}{\mathbb C}
\providecommand{\FF}{\mathbb F}
\providecommand{\NN}{\mathbb N}
\providecommand{\QQ}{\mathbb Q}
\providecommand{\RR}{\mathbb R}
\providecommand{\ZZ}{\mathbb Z}
\providecommand{\dg}{^\circ}
\providecommand{\ii}{\item}
\providecommand{\setdiff}{\smallsetminus}
\providecommand{\alert}[1]{{\sffamily\textbf{\textcolor{blue}{#1}}}}
\providecommand{\inv}{^{-1}}
\providecommand{\mb}[1]{\mathbf{#1}}
\providecommand{\dif}{\;d}

\reversemarginpar

\newenvironment{soln}{\begin{proof}[Solution]}{\end{proof}}
\newenvironment{walkthrough}{\noindent \textbf{Walkthrough.}}{\medskip}
\newlist{walk}{enumerate}{3}
\setlist[walk]{label=\bfseries (\alph*)}

\providecommand{\pow}{\operatorname{Pow}}
\providecommand{\vocab}[1]{{\textbf{\textcolor{ForestGreen}{#1}}}}
\providecommand{\lcm}{\operatorname{lcm}}
\providecommand{\abs}[1]{\left|#1\right|}

\usepackage{asymptote}
\begin{asydef}
	defaultpen(fontsize(10pt));
	unitsize(1cm);
	size(8cm); // set a reasonable default
	usepackage("amsmath");
	usepackage("amssymb");
	settings.tex="pdflatex";
	settings.outformat="pdf";
	// Replacement for olympiad+cse5 which is not standard
	import geometry;
	// recalibrate fill and filldraw for conics
	void filldraw(picture pic = currentpicture, conic g, pen fillpen=defaultpen, pen drawpen=defaultpen)
	{ filldraw(pic, (path) g, fillpen, drawpen); }
	void fill(picture pic = currentpicture, conic g, pen p=defaultpen)
	{ filldraw(pic, (path) g, p); }
	// some geometry
	pair foot(pair P, pair A, pair B) { return foot(triangle(A,B,P).VC); }
	pair orthocenter(pair A, pair B, pair C) { return orthocentercenter(A,B,C); }
	pair centroid(pair A, pair B, pair C) { return (A+B+C)/3; }
	// cse5 abbreviations
	path CP(pair P, pair A) { return circle(P, abs(A-P)); }
	path CR(pair P, real r) { return circle(P, r); }
	pair IP(path p, path q) { return intersectionpoints(p,q)[0]; }
	pair OP(path p, path q) { return intersectionpoints(p,q)[1]; }
	path Line(pair A, pair B, real a=0.6, real b=a) { return (a*(A-B)+A)--(b*(B-A)+B); }
	// cse5 more useful functions
	picture CC() {
		picture p=rotate(0)*currentpicture;
		currentpicture.erase();
		return p;
	}
	pair MP(Label s, pair A, pair B = plain.S, pen p = defaultpen) {
		Label L = s;
		L.s = "$"+s.s+"$";
		label(L, A, B, p);
		return A;
	}
	pair Drawing(Label s = "", pair A, pair B = plain.S, pen p = defaultpen) {
		dot(MP(s, A, B, p), p);
		return A;
	}
	path Drawing(path g, pen p = defaultpen, arrowbar ar = None) {
		draw(g, p, ar);
		return g;
	}
\end{asydef}

\usepackage{mathrsfs}

\ihead{\footnotesize MAT 314 HW02}
\ohead{\footnotesize Last compiled \today}

\title{\vspace{-2em}MAT 314 HW02}
\author{\large Aditya Pahuja}
\date{\large Due 02/20}
\begin{document}
\maketitle

\begin{problem*}[1.1]
	Show that the image of a representation of dimension $1$
	of a finite group is a cyclic group.
\end{problem*}

\begin{soln}
	Let $G$ be a finite group and $\rho\colon G\to \CC^\times$
	a one-dimensional representation (identifying $GL_1(\CC)$ with
	the multiplicative group $\CC^\times = \CC\setdiff\{0\}$).
	We know by Lagrange's theorem that, if $|G| = n$, then for any $g\in G$,
	\[1 = \rho(1) = \rho(g^n) = \rho(g)^n,\]
	so $\rho(g) = \omega^k$ for some integer $k$ with $\omega = e^{2\pi i/n}$. Thus the image of $G$ is a subgroup of
	\[\{\omega^k\colon 0\le k\le n\}\]
	under multiplication, which is a cyclic group,
	and subgroups of cyclic groups are cyclic.
\end{soln}

\begin{problem*}[3.1]
	Let $G$ be a cyclic group of order $3$.
	The matrix $A = \begin{pmatrix}
		-1&-1\\1&0
	\end{pmatrix}$ has order $3$, so it defines a matrix representation of $G$.
	Use the averaging process to produce a $G$-invariant form
	from the standard Hermitian product $X^*Y$ on $\CC^2$.
\end{problem*}

\begin{soln}
	The form generated by the averaging process is
	\[\langle x,y\rangle = \frac1{|G|}\sum_{g\in G} (gx)^*(gy)
	= \frac13\sum_{k=0}^3 x^*(A^k)^*A^ky.\]
	We can compute
	\begin{align*}
		(A^0)^*A^0 &= \begin{pmatrix}
			1&0\\0&1
		\end{pmatrix}\\
		(A^1)^*A^1 &= \begin{pmatrix}
			2&1\\1&1
		\end{pmatrix}\\
		(A^2)^*A^2 &= \begin{pmatrix}
			1&1\\1&2
		\end{pmatrix}
	\end{align*}
	so the average is
	\[\langle x,y\rangle = x^*\begin{pmatrix}
		4/3&2/3\\2/3&4/3
	\end{pmatrix}y,\]
	which, in terms of the components of $x$ and $y$, can be written as
	\[\langle x,y\rangle =
	\boxed{
		\frac{4\ol{x_1}y_1 + 2\ol{x_1}y_2 +2\ol{x_2}y_1 + 4\ol{x_2}y_2}{3}
	}.\qedhere\]
\end{soln}

\begin{problem*}[4.10]
	(a) Find the missing rows in the character table below.
	
	(b) Determine the orders of the elements $a$, $b$, $c$, $d$.
	
	(c) Show that the group $G$ with this character table
	has a subgroup $H$ of order $10$, and describe this subgroup
	as a union of conjugacy classes.
	
	(d) Decide whether $H$ is $C_{10}$ or $D_5$.
	
	(e) Determine all normal subgroups of $G$.
	
	\begin{center}
		\begin{tabular}{c|ccccc}
			&(1)&(4)&(5)&(5)&(5)\\
			&1&$a$&$b$&$c$&$d$\\
			\hline
			$\chi_1$&1&1&1&1&1\\
			$\chi_2$&1&1&$-1$&$-1$&1\\
			$\chi_3$&1&1&$-i$&$i$&$-1$\\
			$\chi_4$&1&1&$i$&$-i$&$-1$
		\end{tabular}
	\end{center}
\end{problem*}

\begin{soln}
	(a) There is one missing row, since $G$ has $5$ conjugacy classes
	hence $5$ representations. Let $\chi_5$ be the missing character.
	\begin{itemize}
		\ii We must have $\chi_5(1) = 4$ because
		\[\chi_1(1)^2 + \chi_2(1)^2 + \chi_3(1)^2 + \chi_4(1)^2 + \chi_5(1)^2
		= |G| = 20.\]
		\ii Since the columns of a character table are orthogonal,
		comparing the first and second columns gives $\chi_5(a) = -1$,
		and comparing the first column to the rest of the columns shows that
		$\chi_5(b) = \chi_5(c) = \chi_5(d) = 0$
		(or one can use the fact that the characters are orthonormal to get
		\begin{align*}
			20 &= |\chi_5(1)|^2 + |\chi_5(a)|^2 + |\chi_5(b)|^2 +
			|\chi_5(c)|^2 + |\chi_5(d)|^2\\
			&= 20 + |\chi_5(b)|^2 + |\chi_5(c)|^2 + |\chi_5(d)|^2\\
			&\ge 20
		\end{align*}
		for the same result because squares of real numbers are nonnegative).
	\end{itemize}
	
	(b) Note that $b$ is the inverse of an element in the conjugacy class of $c$
	because the column of $c$ is the only one conjugate with the column of $b$,
	which means that $b$ and $c$ have the same order.
	Look at $\chi_3$; it is a $1$-dimensional representation
	($\chi_3(1) = 1$), so $\chi_3(G)$ is a cyclic group (Problem 1.1).
	The order of $\chi_3(b)$ in this cyclic group
	divides the order of $b$, meaning the order of $b$ is divisible by $4$.
	This means $b$ either has order $4$ or $20$,
	but the latter would imply that $G$ is cyclic, which is impossible
	because the conjugacy classes in an abelian group are singleton sets,
	which the table disagrees with. Hence \framebox{$b$ and $c$ have order $4$}.
	
	For $a$ and $d$, since one-dimensional representations are homomorphisms
	to $\CC^\times$, we can conclude that the union of
	the conjugacy classes of $1$ and $a$ (kernel of $\chi_3$)
	is a normal subgroup of $G$, as is the union of
	the classes of $1$, $a$, $d$ (kernel of $\chi_2$).
	
	The former subgroup is isomorphic to $\ZZ/5\ZZ$ because
	groups of prime order are cyclic, so \framebox{the order of $a$ is $5$}.
	
	For the latter, subgroup all five elements in the class of $d$
	have the same order.
	This means that the elements' common order,
	and in particular \framebox{the order of $d$, must be 2}.
	The common order can't be $5$ because elements of order $5$
	come in groups of $4$ (it's easy to show that two subgroups of order $5$
	either intersect at the identity or are the same subgroup),
	and the common order can't be $10$ because the existence of
	elements of order $10$ implies the existence of elements of order $2$,
	which couldn't exist here --- we don't have enough conjugacy classes.
	
	(c), (d) We showed above that $H$ must be the group
	consisting of the conjugacy classes of $1$, $a$, and $d$.
	Moreover, we must have $H\cong D_5$ because, as we showed above,
	$G$ has no elements of order $10$.
	
	(e) From the character table, we know that the trivial subgroup,
	$G$, the subgroup of order $10$, and the subgroup of order $5$
	are all normal in $G$. We will show that these are the only one.
	
	We know that normal subgroups are always unions of conjugacy classes,
	so it remains to rule out any other combinations of unions.
	Let $N$ be a normal subgroup of $G$.
	If the conjugacy class of $b$ or $c$ is contained in $N$,
	then so is the other one of the two (because the inverse of
	$b$ is in the class of $c$ and vice versa), and $1\in N$ as well.
	This totals to $11$ elements already, so the only subgroup
	(let alone a normal one) containing these elements is $G$,
	since $20$ is the only divisor of $20$ at least as large as $11$.
	
	Now suppose that the classes of $b$ and $c$ are not subsets of $N$.
	If the conjugacy class of $d$ is contained in $N$, then
	the class of $d$ and $1$ together provide us with $6$ elements,
	so, again noting that $|N|$ divides $|G|$, we get that
	$N$ is the union of the classes of $1$, $a$, and $d$
	or $N = G$, both of which were listed above.
	
	Finally, the only remaining consideration is when
	the conjugacy classes of $b$, $c$, and $d$ are all outside $N$,
	in which case we are either left with the normal subgroup generated by
	$1$ and $a$ or the trivial subgroup, both of which we have counted.
\end{soln}

\begin{problem*}[5.5]
	Prove that the one-dimensional characters of a finite group $G$
	form a group under multiplication of functions.
	This group is called the \vocab{character group} of $G$,
	and is often denoted by $\widehat G$. Prove that if $G$ is abelian,
	$|\widehat G| = |G|$ and $\widehat G\cong G$.
\end{problem*}

\begin{soln}
	Let $\chi,\chi'\in\widehat G$. Then, for any $g,h\in G$,
	\[\chi(gh)\cdot \chi'(gh) = \chi(g)\chi(h)\cdot\chi'(g)\chi'(h)
	= (\chi\chi')(g)\cdot (\chi\chi')(h),\]
	which proves that the product of two characters is a character,
	hence $\widehat G$ is closed under multiplication.
	Multiplication is also an associative operation,
	and there is an identity $\chi_1\equiv 1$ and an inverse
	$\chi\inv = \frac{1}{\chi}$ (noting that the codomain of $\chi$
	is the multiplicative group of complex numbers, so in particular
	$\chi(g)\ne 0$ for all $g$), so $\widehat G$ is indeed a group.
	
	To show that $G$ is isomorphic to its character group if it is abelian,
	we first prove that this statement is true when $G\cong\ZZ/n\ZZ$
	is cyclic. In this case, the characters (of which there are $n = |G|$,
	because $G$ has $n$ conjugacy classes) are uniquely determined by
	the image of a generator $g$ of $G$. We will assign a character to
	each element $g^k$ of $G$ by
	\[\chi_{g^k}(g) = \omega^k\]
	where $\omega = \exp(\frac{2\pi i}n)$ is a primitive $n$th root of unity.
	Then, we can see that for any characters
	$\chi_{g^a},\chi_{g^b}\in\widehat G$,
	\[\chi_{g^a}(g)\cdot \chi_{g^b}(g) =
	\omega^a\cdot\omega^b = \omega^{a+b} = \chi_{g^{a+b}}(g),\]
	which proves in general that $\chi_{g^a}\chi_{g^b} = \chi_{g^{a+b}}$
	because $g$ generates $G$, so $G\cong\widehat G$.
	
	Now, we show that we can patch together cyclic groups
	to get at arbitrary finite abelian groups.
	\begin{claim*}
		Suppose that finite abelian groups $G$, $H$ are
		isomorphic to their character groups. Then,
		\[\widehat G\times\widehat H\cong \widehat{G\times H}.\]
	\end{claim*}
	\begin{proof}
		Let $g_1 = 1_G$, $g_2$, \dots, $g_n$ be the $n$ elements of $G$
		and $h_1 = 1_H$, $h_2$, \dots, $h_m$ be the $m$ elements of $H$.
		Then, given an element $(\chi_{g_i}, \chi_{h_j})\in
		\widehat G\times\widehat H$, we assign to $(g_i,h_j)\in G\times H$
		the character of $G\times H$ defined by
		\[\chi_{(g_i,h_j)}(g,h) := \chi_{g_i}(g)\cdot \chi_{h_j}(h)\]
		for any $(g,h)\in G\times H$.
		To see that this is indeed a character, we
		just leverage the fact that $\chi_{g_i}$ and $\chi_{h_j}$
		are characters of $G$ and $H$ to conclude that
		\begin{align*}
			\chi_{(g_i,h_j)}((g,h)\cdot(g',h'))
			&=\chi_{(g_i,h_j)}(gg',hh')\\
			&= \chi_{g_i}(gg')\cdot\chi_{h_j}(hh')\\
			&= \chi_{g_i}(g)\chi_{h_j}(h)\cdot\chi_{g_i}(g)\chi_{h_j}(h)\\
			&= \chi_{(g_i,h_j)}(g,h)\cdot\chi_{(g_i,h_j)}(g',h')
		\end{align*}
		so that $\chi_{(g_i,h_j)}$ is a homomorphism $G\times H\to\CC$.
		Moreover, $\chi_{(g_i,h_j)}$ is uniquely determined by
		the images of $G\times\{1_H\}$ and $\{1_G\}\times H$,
		which are in turn in bijection with the images of
		$\chi_{g_i}$ and $\chi_{h_j}$, so we have
		exactly $mn = |G\times H|$ characters
		and our map $(\chi_{g_i}, \chi_{h_j})\mapsto\chi_{(g_i,h_j)}$
		is an injective homomorphism between finite groups,
		hence it is an isomorphism.
		Thus we have exhibited an isomorphism between
		$\widehat G\times\widehat H$ and $\widehat{G\times H}$, as desired.
	\end{proof}
	
	To finish, we recall that every finite abelian group $G$
	is a direct product of finite cyclic groups,
	so the claim decrees that $G$ must be isomorphic to $\widehat G$.
\end{soln}
\pagebreak
\begin{problem*}[7.1]
	Prove a converse to Schur's lemma: If $\rho$ is a representation,
	and if the only $G$-invariant linear operators on $V$
	are multiplications by scalars, then $\rho$ is irreducible.
\end{problem*}

\begin{soln}
	Suppose $\rho$ were reducible. Then, by Maschke's theorem,
	we can write $V$ as the direct sum of two $G$-invariant subspaces
	$U$, $W$. Consider the map $\pi\colon V\to W$ given by
	$\pi(v) = u$ where $v = u + w$ is the unique representation of $V$
	as the sum of a vector $u\in U$ and a vector $w\in W$. We get
	\[\pi(gv) = \pi(gu + gw) = \pi(gu) + \pi(gw) = \pi(gw)
	= gw = g\pi(w) = g\pi(v),\]
	using the fact that $gu\in U$, $gw\in W$, and $\pi(a) = a$ for all $a\in W$.
	This is therefore a non-scalar $G$-invariant linear operator on $G$,
	as desired.
\end{soln}




















\end{document}